%!TeX root=../carlo_master.tex
\chapter{Conclusion}\label{sec:conclusion}

This dissertation has argued that the decision stage of cost-sensitive classification is a learnable component, separable from the inference estimator that supplies it and small enough to be fit reliably from a held-out partition. The present chapter closes the work in four parts: a chapter-by-chapter summary of the contributions, a practitioner-facing recipe extracted from the campaigns, the scope decisions that bound the claims, and the follow-up directions they invite.

\section{Summary of contributions}

This dissertation treats the decision stage of cost-sensitive classification as a separately optimizable component, distinct from the inference stage that produces the probability estimates it consumes. The motivation is that the analytical plug-in rule incurs excess risk when the probability estimator is miscalibrated under an asymmetric cost structure, the regime that modern classifiers routinely occupy. A held-out optimization of the decision layer corrects decision-level miscalibrations that the analytical rule cannot.

Chapter~\ref{sec:theory} gives the structural reason this correction is admissible: a complexity asymmetry between the two stages. The inference model may be high-complexity and miscalibrated, while the decision class has dimension polynomial in $K$. On a validation sample of size $m$ drawn independently of the estimator fit, the decision layer therefore admits a generalization bound of order $\sqrt{d_\mathrm{dec}/m}$ (Proposition~\ref{prop:gen_bound}).

This capacity argument is necessary, not sufficient. It establishes that the decision layer is learnable in principle, but not that the gain materializes on a given dataset. On a single small held-out split the empirical minimizer can underperform the analytical plug-in. Across both campaigns the improvement is reliable only when the decision layer is selected by cross-validation over the full training pool. The capacity bound makes the procedure statistically sound; it is the cross-validated estimation, not the bound, that produces the improvement in practice.

In the binary setting (Chapters~\ref{sec:binary} and~\ref{sec:binary_experiments}) we derive the analytical Elkan threshold and characterize its fragility under probability miscalibration. The campaign spans fourteen datasets, ten base classifiers, and ten objectives. EmpiricalCV wins or ties the per-row best on $87\%$ of cost cells and $99.9\%$ of maximization-metric cells.

In the multiclass setting (Chapters~\ref{sec:multiclass} through~\ref{sec:multiclass_experiments}) we develop a taxonomy of partition-optimization strategies (the parallel shift and linear class scaling of \citet{obrien2008costsensitive}, the probability-ratio thresholding of \citet{lachiche2003improving} and the multidimensional thresholding of \citet{marchetti2025multidimensional}, each equivalent under $0/1$ loss to one of the O'Brien parameterizations, and a priority-based sequential cascade), grouping the matrix-parameterized classes $\mathcal{H}_{PS} \cup \mathcal{H}_{Sc} \cup \mathcal{H}_{Full}$ under a single framework and identifying $\mathcal{H}_{PS} \cup \mathcal{H}_{Sc}$ as the $K-1$-parameter vector-tuned subset. We compute Natarajan-dimension bounds for the convex-partition score classes: $d_{Nat}(\mathcal{H}_{PS}) = K-1$, $d_{Nat}(\mathcal{H}_{Sc}) \leq K-1$, and $d_{Nat}(\mathcal{H}_{Full}) \leq K^2$. The unconstrained Voronoi class $\mathcal{H}_{Vor}$ of Chapter~\ref{sec:voronoi_tessellation} properly contains $\mathcal{H}_{PS}$ (Proposition~\ref{prop:ps_in_vor}), with strict population risk separation from $\mathcal{H}_{PS}$ exhibited on a $K{=}3$ witness under symmetric $0/1$ loss via the DLC subclass $\mathcal{H}_{Vor}^{\mathrm{DLC}}$ (Theorem~\ref{thm:rotational_dominance}). The class actually fit is $\mathcal{H}_{Vor}^{\mathrm{DLC}} \subseteq \mathcal{H}_{Vor}^{\Delta}$ at $d_{Nat} \leq K^2$, rather than the $d_{Nat} \leq K^2 + K$ unconstrained class on which Proposition~\ref{prop:ps_in_vor} relies.

The cross-validated tuning variants lead the multiclass per-row comparison as in the binary setting. On cost objectives three variants are tied at the top within sampling noise (FullMatrixPartitioningCV at $80\%$, RiskPartitioningCV at $79\%$, LinearScalingCV at $78\%$ win-or-tie); on maximization metrics ThresholdRatioCV leads at $90\%$. The richer Voronoi class does not improve on these: counting cells where a Voronoi variant is the unique SE-overlap winning family, the rate is $0.1\%$ on cost cells and $1.2\%$ on maximization cells. This is consistent with Proposition~\ref{prop:ps_in_vor}, which guarantees only no-worse population risk; the strict advantage of Theorem~\ref{thm:rotational_dominance} requires a rotational regime the OpenML benchmark does not appear to concentrate in.

Chapter~\ref{sec:misspec_sensitivity_K} adds a second-order observation about how the decision-layer advantage depends on $K$. The lead of the best strategy over the bare analytical rule on cost cells narrows from $18$ points in binary ($87\%$ versus $69\%$ win-or-tie) to $11$ points in multiclass ($80\%$ versus $69\%$), so the room for a decision-layer correction narrows as $K$ grows. The dimensional-dilution mechanism predicts this direction qualitatively: a diffuse misspecification of $\boldsymbol{\eta}$ of fixed Euclidean radius projects onto a scalar cost functional with expected squared overlap of order $1/(K-1)$ (Lemma~\ref{lem:dimensional_dilution}), verified to that exponent in a synthetic setting where its assumptions hold by construction. The benchmark comparison and the synthetic mechanism are not equivalent: the benchmark mixes the geometric effect with sample-complexity effects on the tuning side, and the lemma's quantitative exponent does not transfer at finite perturbation magnitude in the variance-scaling and mean-shift experiments.

\section{A practitioner recipe}\label{sec:practitioner_recipe}

The campaigns support a short rule of thumb. Use a cross-validated decision-layer tuner whenever the training data supports it: EmpiricalCV in binary and, in multiclass, FullMatrixPartitioningCV, RiskPartitioningCV, or LinearScalingCV for cost-decomposable objectives, and ThresholdRatioCV for non-decomposable maximization metrics. These strategies lead every campaign aggregate in this dissertation, from $87\%$ and $99.9\%$ for EmpiricalCV on binary cost and maximization cells to the $78$--$80\%$ band on multiclass cost cells and $90\%$ on multiclass maximization cells. The one exception is a well-calibrated estimator (the three deep tabular models in their supported regimes) paired with a cost-decomposable objective, where the bare Analytical plug-in is competitive within sampling noise and consumes no held-out data.

Two operational caveats follow. First, the empirical family beats the analytical plug-in only when the full training pool funds the tuning: single-split Empirical on the $10\%$ validation partition wins or ties just $42\%$ of binary cost cells, below the bare Analytical's $69\%$, so an in-budget held-out split is not a substitute for cross-validating over the full $80\%$ pool. Second, the richer partition classes earn their parameters only when the data supplies a matching signal. VoronoiTuning is appropriate when the estimator's miscalibration rotates the optimal boundaries (the regime constructed in Section~\ref{sec:boundary_rotation_experiment}); the OpenML pool does not concentrate there, and the vector-tuned and full-matrix CV variants meet or beat it within sampling noise on both objective groups.

\section{Limitations}

Nine scope decisions bound the claims above and invite extension.

\begin{enumerate}
    \item The evaluation protocol is fixed throughout the campaigns: $80/20$ train-test, with a $70/10$ train/validation split inside the train-full pool for the single-split tuners and $5$-fold stratified cross-validation over the full $80\%$ pool for the CV variants. How the per-strategy rankings vary with these knobs (a tighter or looser validation partition, a different fold count, or repeated CV) is not investigated.
    \item The per-strategy search ranges (for instance $\boldsymbol{\Delta} \in [-1, 1]^{K \times K}$ for FullMatrixPartitioning, $b_k \in [-1, 1]$ for MultidimensionalThreshold, $\tau_k \in [0.01, 1.5]$ for ThresholdRatio) are fixed by a single up-front choice rather than adapted per dataset, classifier, objective, or cost-matrix scale. A more carefully designed search space (rescaling $\boldsymbol{\Delta}$ by $C_{\max}$ to track the anchor magnitude, widening the parallel-shift band on objectives where the best $\mathbf{b}$ sits near the search boundary, or per-regime tuning) could move any of the tuning variants in either direction, so the per-strategy rankings reflect both the strategies themselves and the specific ranges adopted here.
    \item The Voronoi optimization relies on stochastic local search, since the empirical risk is piecewise constant in the prototype configuration. A differentiable surrogate (for instance, a softmax-temperature relaxation of the nearest-prototype assignment) would admit first-order methods. The sampling distribution that feeds the search is a separate axis: the experiments use the sequential $d^2$ sampler of Section~\ref{sec:voronoi_samplers}, adopted as the only one of the six surveyed there that satisfies~\eqref{eq:dlc} with probability one across the full $K$ sweep, but a sampler biased toward the analytical-parallel-shift seed, or one that relaxes the DLC in exchange for wider coverage of $\mathcal{H}_{Vor}^{\Delta}$, is unexplored. Either change may reduce the gap between the matrix-parameterized methods and Voronoi observed on the OpenML benchmark.
    \item The multiclass benchmark covers $K \leq 10$. Whether the empirical Pareto-dominance of the matrix-parameterized methods over Voronoi inverts at larger $K$, where the additional rotational degrees of freedom grow quadratically in $K$, is left open.
    \item The models are evaluated at default parameters only. A controlled study of how decision-layer tuning interacts with base-estimator calibration quality across model capacities would refine the practical recommendations of Section~\ref{sec:practitioner_recipe}.
    \item The benchmark draws exclusively from tabular OpenML datasets and tabular-oriented base classifiers (logistic regression, gradient-boosted trees, random forests, and three deep tabular architectures). Whether the per-strategy rankings and the recommendations of Section~\ref{sec:practitioner_recipe} transfer to non-tabular multiclass settings (image classification with convolutional or transformer backbones, text classification with language-model embeddings, time-series, audio, or graph classification) is an open question; the calibration profiles and typical cost-asymmetry regimes practitioners face in those domains may differ enough that a separate benchmark is required.
    \item The calibrate-then-decide strategies fit on the $10\%$ validation partition while their direct-tuning CV counterparts use the full $80\%$ pool, a roughly sevenfold gap in label budget. The natural matched-budget protocol is to cross-validate the calibrator parameters as hyperparameters, the same treatment the CV tuning variants give to thresholds and partition parameters; this is left outside scope, and the calibrate-then-decide rankings are read as conditional on the budget asymmetry.
    \item The $1/(K-1)$ dilution exponent is a property of the linearized cost functional under an isotropic perturbation of the simplex tangent space. The variance-scaling and mean-shift experiments reproduce the K-monotone ordering but not the exact exponent at finite perturbation magnitude. Whether realistic estimator misspecification on a real dataset is sufficiently diffuse for this mechanism to dominate over the worst-case binary-subproblem regime is an open question.
    \item Running the campaigns end-to-end is non-trivially expensive (Appendix~\ref{app:compute} documents the wall-clock and cloud-spend footprint, together with the harness optimisations that brought the multiclass campaign down by an order of magnitude). The compute floor constrained how many of the axes flagged in the limitations above could be explored: sweeping the validation-split size, the fold count, or the per-strategy search ranges, or adding seeds, datasets, or estimator-capacity ablations, would have required compute beyond what was budgeted for this work.
\end{enumerate}

\section{Future work}

The two-stage formulation leaves four natural follow-up directions.

\paragraph{Cost-into-loss training.} The alternative paradigm set aside in this dissertation folds the cost structure into the training objective rather than pushing it to the decision stage. Cost-weighted cross-entropy, focal loss, and differentiable surrogates for confusion-matrix metrics, including the Score-Oriented Loss framework of \citet{marchetti2022sol}, are the main candidates. The two routes are not mutually exclusive: a cost-weighted training objective can still be paired with a tuned decision rule. A controlled comparison on the benchmarks of Chapters~\ref{sec:binary_experiments} and~\ref{sec:multiclass_experiments} would quantify when integrating the cost into the loss replaces, complements, or is dominated by tuning at the decision stage.

\paragraph{Richer cost structure.} The cost matrix is treated here as a fixed constant. Many applied settings carry costs that depend on the input (fraud detection where the loss scales with transaction value, medical triage where the misclassification cost depends on patient characteristics), on time (drift, deteriorating signal, deadlines), or on group membership (fairness constraints expressed as group-conditional miscalibration penalties). Extending the partition-optimization framework to $\mathbf{C}(x)$, $\mathbf{C}(t)$, or $\mathbf{C}(g)$ requires reformulating $\mathcal{H}_{PS}$, for instance, as parameter fields over the conditioning variable rather than constants on the simplex; the analytical plug-in rule extends in form, since it is already a pointwise minimization $\arg\min_k \mathbf{c}_k^\top \boldsymbol{\eta}(x)$ at each $x$ and only the cost vector needs to be queried at the conditioning variable, but the empirical-tuning machinery does not transfer immediately, since the partition parameters themselves become $x$-, $t$-, or $g$-indexed.

\paragraph{Argmax as the default elsewhere in ML.} The decision-stage perspective developed here is one specialization of a recurring pattern: in many machine-learning pipelines a selection step (whether final or repeated inside the pipeline) is a learnable design choice that the standard implementation forgoes by defaulting to argmax or its top-$k$ cousin. In autoregressive language-model decoding, greedy next-token selection gives way to temperature sampling, nucleus and top-$k$ sampling \citep{holtzman2020curious}, beam search, and contrastive or speculative decoding. In reinforcement learning and contextual bandits \citep{sutton2018reinforcement}, the entire exploration literature exists to improve on greedy argmax over value estimates ($\varepsilon$-greedy, Boltzmann sampling, UCB, Thompson sampling); modern actor-critic methods (PPO, SAC) replace deterministic argmax policies with stochastic ones, with entropy regularization explicitly penalizing argmax collapse, and risk-sensitive RL replaces expected-return-argmax with CVaR or distortion-based action rules \citep{tamar2015cvar}, a direct analogue of the cost-sensitive theme here. The pattern recurs in semantic segmentation (per-pixel argmax versus learned decoders or CRF post-processing \citep{krahenbuhl2011efficient}), in retrieval and recommendation (top-$k$ cutoffs and ranking-aware thresholds), in selective classification (predict-vs-abstain thresholds, \citealp{geifman2017selective}), and in object detection (the confidence threshold of non-maximum suppression, \citealp{bodla2017soft}). Each of these has a learnable or cost-aware variant that the standard pipeline silently forgoes.

\paragraph{Releasing the library.} The decision strategies of this dissertation are implemented in \texttt{threshold\_opt}, a Python library that follows the scikit-learn \citep{pedregosa2011scikit} estimator conventions: every strategy inherits from a common \texttt{DecisionStrategy} base class and exposes a \texttt{fit(X, y, objective)} / \texttt{predict(X)} pair, with composable post-hoc calibrators, a vectorised cost-aware metric implementation, and the harness optimisations of Appendix~\ref{app:compute} on the experiment side. Finalising the public-API documentation, type annotations, and PyPI release is a concrete follow-up. Decision-stage tuning is a step that practitioners commonly skip, defaulting to the cost-blind argmax; the campaigns in this dissertation suggest that the gap on miscalibrated estimators and non-decomposable objectives is large enough to warrant a library that makes adopting the decision-layer step as routine as fitting a calibrator.

The argmax that ships with every off-the-shelf classifier is a choice, not a constant. When the estimator is well-calibrated and the objective is cost-decomposable, Bayesian decision theory supplies the answer in closed form: the analytical Elkan threshold and its multiclass partition boundary are competitive without any held-out tuning. When the estimator is miscalibrated or the objective is non-decomposable, the answer is a learnable rule, fit on a validation partition with the same care as a calibrator. The decision stage is small enough to learn; the cost of leaving it at the cost-blind default is large enough to matter.
